\section{Limitations}
\label{sec:limitations}

\begin{itemize}
\item \textbf{Oracle-assisted recovery.} The recovery loop is gated by an
oracle, so every reported recovery is oracle-assisted rather than
autonomous self-correction. Stopping at the first PASS hides later
revisions that could undo a correct answer, and an evaluator that can be
queried repeatedly favors the team that queries it most.

\item \textbf{Scope and exploratory status.} We test one \PER{}
organization on text-based reasoning tasks, with an evaluator that is
itself one of the models under study. The main comparison and the
forced-submission run are exploratory, added after the freeze. The
forced-submission run is a single run for one actor on one domain, and it
removes the whole gate, not only the clause forbidding approval after a
failure.

\item \textbf{Feedback and cost.} Equal access to feedback does not
equalize the realized number or content of verdicts, and the two teams see
the feedback in different places. We do not show that any submission
policy is cost-optimal (further disclosures in Appendix~\ref{app:details}).

\item \textbf{Scoring and logging defects.} The runner graded the last
\emph{submitted} answer but recorded the last \emph{held} one, the team's
current answer. For gemma-4 on Omni-MATH they disagree on $17\%$ of
role-structured first-pass-wrong trajectories. Recoding them narrows the
recovery gap and moves the verdict-matched difference toward roles
(\S\ref{subsec:mechanism}, Table~\ref{tab:robust}), but the primary
analysis keeps the pre-specified rule. A logging defect hides most of
gpt-oss's role-free reasoning traces from the record, not from the agents,
and does not create the collapse (Appendix~\ref{app:results}).
\end{itemize}
